Simultaneous Localization and Mapping (SLAM) is a foundational capability
for autonomous robots, augmented reality (AR), and autonomous navigation.
Traditional SLAM systems---whether filter-based~\cite{qin2018vinsmono,
xu2022fastlio2} or graph-based~\cite{campos2021orbslam3}---produce sparse
or semi-dense maps that suffice for localization but lack the photo-realistic
detail required for immersive visualization, semantic understanding, or
high-fidelity simulation.  The introduction of Neural Radiance Fields
(NeRF)~\cite{mildenhall2020nerf} and, more critically, 3D Gaussian
Splatting~\cite{kerbl2023_3dgs} has catalyzed a paradigm shift: SLAM systems
now aim to simultaneously estimate the camera trajectory and reconstruct a
\emph{photo-realistic, dense, and renderable} 3D map of the environment.

3DGS represents scenes as sets of anisotropic 3D Gaussian primitives with
position, covariance, opacity, and spherical-harmonic color coefficients.
Unlike NeRF's costly volumetric ray-marching, 3DGS employs fast
tile-based rasterization, achieving real-time rendering at over 100~FPS on
consumer GPUs.  This efficiency, combined with differentiable optimization,
makes 3DGS the ideal map representation for online SLAM.

However, the quality of a Gaussian-based SLAM map is fundamentally limited by
the \emph{depth source} that seeds and supervises the Gaussians.  Early 3DGS
SLAM systems~\cite{yan2024gsslam,keetha2024_splatam,matsuki2024monogs}
operated with RGB-D cameras indoors, where dense active depth is readily
available.  When deployed in large-scale outdoor environments, these systems
degrade under violent ego-motion, varying illumination, textureless regions,
and---most critically---the limited range of consumer RGB-D sensors.

Three complementary strategies have emerged to provide the geometric priors
needed for robust Gaussian SLAM beyond the indoor RGB-D setting:

\begin{enumerate}
    \item \textbf{LiDAR + RGB + IMU fusion} (Section~\ref{sec:lidar}):
    Mechanically or solid-state LiDAR provides precise, long-range, but
    sparse metric geometry.  The camera contributes rich texture, and the
    IMU bridges motion gaps.  This configuration, well-established in
    classical LIO/LIVO
    systems~\cite{xu2022fastlio2,zheng2024_fastlivo2,lin2024_r3live},
    has been extended with Gaussian Splatting mapping by a rapidly growing
    family of works~\cite{lang2026_gaussianlic2,lang2025_gaussianlic,
    hong2025_gslivo,xie2024_gslivm,zhao2025_lvigs,li2026_rmgs,
    tak2026_realtimelidar,he2026lvgsslam,shi2026_litgs}.

    \item \textbf{iToF Depth + RGB + IMU fusion} (Section~\ref{sec:itof}):
    Indirect Time-of-Flight (iToF) sensors offer low-power, compact active
    depth at the cost of extreme sparsity (as few as $8\times8$ points per
    frame) and environmental noise sensitivity.  Recent
    work~\cite{conti2025_tofsplatting} demonstrates that even such sparse
    depth, when fused with multi-view geometry and monocular cues within a
    3DGS framework, can achieve competitive SLAM performance---a crucial
    step toward mobile and AR/VR deployment.

    \item \textbf{Monocular Depth + RGB + IMU fusion} (Section~\ref{sec:monodepth}):
    When no active depth sensor is available, learned monocular depth
    priors~\cite{yang2024depthanythingv2} or feed-forward 3D foundation
    models~\cite{wang2025_vggt} can substitute.  The IMU provides metric
    scale and motion continuity.  This axis spans pure visual-inertial
    Gaussian SLAM~\cite{zhu2026_vigssslam,zhang2026_hislam2,
    seo2026_gaussianflow,zhou2026_geogs,zheng2025_wildgssslam,
    han2026_vggtgssslam,zhang2026_vimgs,deambrogi2026_cgssslam,
    sun2024_mm3dgs,zhu2026_mygosplat,kruzhkov2026_monoemgs} and is the
    most hardware-accessible but algorithmically challenging configuration.
\end{enumerate}

\textbf{Scope and methodology.}  This survey covers papers published or
posted as preprints through October~2026, identified through systematic web
searches using the Doubao search API with keywords combining ``Gaussian
Splatting,'' ``SLAM,'' ``LiDAR,'' ``ToF,'' ``monocular depth,'' ``visual
inertial,'' and related terms.  For each paper, we downloaded the arXiv PDF
(when available), read the full method section, and searched for open-source
code repositories.  Papers already present in the project's related-work
bibliography (130+ entries accumulated over 30+ survey addenda) were
cross-referenced to avoid duplication.  We focus on systems that perform
\emph{online} or \emph{real-time} SLAM with 3DGS as the primary map
representation, excluding offline 3DGS reconstruction and non-SLAM novel-view
synthesis.

\textbf{Contributions.}  This survey provides:
\begin{itemize}
    \item A structured taxonomy of multimodal Gaussian Splatting SLAM along
    three depth-source axes (LiDAR, iToF, monocular depth).
    \item Per-paper analysis of architecture, depth fusion strategy, code
    availability, and key results for 25+ representative systems.
    \item A cross-category comparison identifying the accuracy--cost--robustness
    trade-off landscape.
    \item A discussion of open challenges: depth completion under degeneracy,
    scale consistency, memory-bounded mapping, and real-time deployment on
    embedded platforms.
\end{itemize}
